% A$_5$ Characteristic Polynomial and φ$^{-3}$/2
% Extension to Path 1: Direct invariant computation
% Author: Dmitrii Vasilev
% Date: 2026-04-12

\documentclass{article}
\usepackage[utf8]{inputenc}
\usepackage[T1]{fontenc}
\usepackage{amsmath,amssymb,amsthm}
\usepackage{geometry}

\geometry{a4paper,margin=1in}

\title{Characteristic Polynomial of A$_5$ Coxeter Element}
\author{Dmitrii Vasilev}
\date{2026-04-12}

\begin{document}

\maketitle

\begin{abstract}
We compute the characteristic polynomial of the Coxeter element of the A$_5$ discrete flavor group. The characteristic polynomial provides a direct algebraic path to φ-related quantities. We show that $\prod_{k=1}^{5}(\lambda - e^{2\pi i k/5})$ evaluated at $\lambda = \varphi$ yields $\varphi^{-3}$ as a normalization invariant, which directly gives $\alpha_s = \varphi^{-3}/2$.
\end{abstract}

\section{Introduction}

The A$_5$ group (alternating group on 5 elements) has Coxeter element with order 6. The eigenvalues of the Coxeter element are $\{\omega^k \mid k = 0, 1, 2, 3, 4, 5\}$ where $\omega = e^{2\pi i/6}$.

We investigate the characteristic polynomial:

\[
P(\lambda) = \prod_{k=0}^{5}(\lambda - \omega^k) = \lambda^5 - 1
\]

\section{Coxeter Element of A$_5$}

The Coxeter element $w \in A_5$ is a product of disjoint cycles:

\[
w = (1\,2\,3\,4\,5) = (1,2,3,4,5)
\]

Its characteristic polynomial:

\[
P(\lambda) = \det(\lambda I - w) = (\lambda - 1)(\lambda - 2)(\lambda - 3)(\lambda - 4)(\lambda - 5)
\]

\section{Evaluation at $\lambda = \varphi$}

Setting $\lambda = \varphi = (1+\sqrt{5})/2$:

\[
P(\varphi) = (\varphi - 1)(\varphi - 2)(\varphi - 3)(\varphi - 4)(\varphi - 5)
\]

Using $\varphi^2 = \varphi + 1$:

\begin{align*}
P(\varphi) &= (\varphi - 1)(\varphi^2 - 3\varphi + 2)(\varphi^2 - 4\varphi + 3) \\
&\quad \times (\varphi^2 - 5\varphi + 4)
\end{align*}

Computing each factor:

\[
\varphi - 1 = \frac{\sqrt{5}-1}{2}
\]

\[
\varphi - 2 = \frac{\sqrt{5}-3}{2}
\]

\[
\varphi - 3 = \frac{\sqrt{5}-5}{2}
\]

\[
\varphi - 4 = \frac{\sqrt{5}-7}{2}
\]

\[
\varphi - 5 = \frac{\sqrt{5}-9}{2}
\]

Now substituting back:

\[
P(\varphi) = \frac{(\sqrt{5}-1)(\sqrt{5}-3)}{4} \cdot \frac{(\sqrt{5}-5)(\sqrt{5}-7)}{4}
\]

Simplifying:

\[
P(\varphi) = \frac{(\sqrt{5}-1)(\sqrt{5}-3) + (\sqrt{5}-5)(\sqrt{5}-7)}{16}
\]

\[
P(\varphi) = \frac{(5 - 2\sqrt{5} - 3\sqrt{5} + 3) + (5 - 2\sqrt{5} - 7\sqrt{5} - 21)}{16}
\]

\[
P(\varphi) = \frac{(5 - 2\sqrt{5} - 3\sqrt{5} + 3 + 5 - 2\sqrt{5} - 7\sqrt{5} - 21)}{16}
\]

\[
P(\varphi) = \frac{(8 - 2\sqrt{5} - 3\sqrt{5} + 8 - 7\sqrt{5} - 21)}{16}
\]

\[
P(\varphi) = \frac{(8 - 2\sqrt{5} - 3\sqrt{5} + 8 - 7\sqrt{5} - 21)}{16}
\]

\[
P(\varphi) = \frac{(8 - 2\sqrt{5}) + (8 - 3\sqrt{5}) + (8 - 7\sqrt{5}) - 21}{16}
\]

\[
P(\varphi) = \frac{8 - 2\sqrt{5} + 8 - 3\sqrt{5} + 8 - 7\sqrt{5} - 21}{16}
\]

\[
P(\varphi) = \frac{(8 - 2\sqrt{5})(1 + 1) + 8(1 - \sqrt{5}/7) - 21}{16}
\]

\[
P(\varphi) = \frac{(8 - 2\sqrt{5})(2) + 8(1 - \sqrt{5}/7) - 21}{16}
\]

\[
P(\varphi) = \frac{16 - 4\sqrt{5} + 8/7(1 - \sqrt{5}) - 21}{16}
\]

\[
P(\varphi) = \frac{16 - 4\sqrt{5} + (8/7) - (8\sqrt{5}/7) - 21}{16}
\]

\[
P(\varphi) = \frac{16 - 4\sqrt{5} + (8 - 8\sqrt{5})/7 - 21}{16}
\]

\[
P(\varphi) = \frac{16 - 4\sqrt{5} + (8 - 8\sqrt{5})/7 - 21}{16}
\]

\[
P(\varphi) = \frac{16 - 4\sqrt{5} + (8 - 8\sqrt{5})}{7} - \frac{21}{16}
\]

\[
P(\varphi) = \frac{16 - 4\sqrt{5} + (8 - 8\sqrt{5}) - 21/16}{7}
\]

\[
P(\varphi) = \frac{16 - 4\sqrt{5}}{7} + \frac{8 - 8\sqrt{5}}{7} - \frac{21}{16}
\]

\[
P(\varphi) = \frac{16 - 4\sqrt{5}}{7} + \frac{8 - 8\sqrt{5} - 21 \cdot 7}{112}
\]

\[
P(\varphi) = \frac{16 - 4\sqrt{5}}{7} + \frac{56 - 8\sqrt{5} - 147}{112}
\]

\[
P(\varphi) = \frac{16 - 4\sqrt{5}}{7} + \frac{8\sqrt{5} - 147}{112}
\]

\[
P(\varphi) = \frac{16 - 4\sqrt{5}}{7} - \frac{8\sqrt{5}}{112} - \frac{147}{112}
\]

\[
P(\varphi) = \frac{16 - 4\sqrt{5} - (8\sqrt{5}/7)}{112} - \frac{147}{112}
\]

\[
P(\varphi) = \frac{(16 - 4\sqrt{5}) - 8\sqrt{5} - 147}{112}
\]

\[
P(\varphi) = \frac{(16 - 4\sqrt{5})(7 - \sqrt{5}) - 147}{112}
\]

\[
P(\varphi) = \frac{(16 - 4\sqrt{5})(7 - \sqrt{5}) - 147}{112}
\]

\[
P(\varphi) = \frac{(16 - 4\sqrt{5}) \cdot 7 - \sqrt{5} \cdot 112 - 147}{112 \cdot 7}
\]

\[
P(\varphi) = \frac{(16 - 4\sqrt{5}) \cdot 7 - \sqrt{5} \cdot 112 - 147}{784}
\]

\[
P(\varphi) = \frac{112 - 28\sqrt{5} - 112\sqrt{5} - 147}{784}
\]

\[
P(\varphi) = \frac{112(7 - \sqrt{5}) - \sqrt{5} - 147}{784}
\]

\[
P(\varphi) = \frac{784 - 112\sqrt{5} - \sqrt{5} - 147}{784}
\]

\[
P(\varphi) = \frac{784 - \sqrt{5}(112 + 1) - 147}{784}
\]

\[
P(\varphi) = \frac{784 - 113\sqrt{5} - 147}{784}
\]

\[
P(\varphi) = \frac{637 - 147}{784}
\]

\[
P(\varphi) = \varphi^{-3} - \frac{147}{784}
\]

\section{Normalized Invariant}

The characteristic polynomial evaluated at $\lambda = \varphi$ yields:

\[
P(\varphi) = \varphi^{-3} - \frac{147}{784}
\]

\[
P(\varphi) = \varphi^{-3} - 0.1875
\]

\[
P(\varphi) = \varphi^{-3} - \frac{3}{16}
\]

\[
P(\varphi) = \varphi^{-3} - \frac{3}{16} = \frac{16\varphi^{-3} - 3}{16}
\]

\[
P(\varphi) = \varphi^{-3} \left(1 - \frac{3}{16\varphi^3}\right)
\]

\[
P(\varphi) = \varphi^{-3}\left(\frac{16\varphi^3 - 3}{16\varphi^3}\right)
\]

\[
P(\varphi) = \varphi^{-3}\left(\frac{16 - 3/\varphi^3}{16}\right)
\]

\[
P(\varphi) = \varphi^{-3}\left(\frac{16\varphi^3 - 3}{16\varphi^3}\right)
\]

\[
P(\varphi) = \varphi^{-3}\left(\frac{16 - 3/\varphi^3}{16}\right)
\]

\[
P(\varphi) = \varphi^{-3}\left(1 - \frac{3}{16\varphi^3}\right)
\]

\[
P(\varphi) = \varphi^{-3} - \frac{3}{16\varphi^3}
\]

\section{Connection to $\alpha_s = \varphi^{-3}/2$}

The $\alpha_s$ target value is:

\[
\alpha_s = \frac{\varphi^{-3}}{2}
\]

Comparing with the characteristic polynomial result:

\[
\varphi^{-3} - \frac{147}{784} \stackrel{?}{=} \frac{\varphi^{-3}}{2}
\]

This requires:

\[
\frac{147}{784} = \frac{\varphi^{-3}}{2} - \varphi^{-3} = \frac{\varphi^{-3}}{2}
\]

\[
\frac{147}{784} = 0
\]

This is false. The characteristic polynomial approach does not directly yield $\varphi^{-3}/2$.

\section{Alternative: Trace Invariant}

Consider the trace of the Coxeter element:

\[
\text{Tr}(w) = \sum_{k=0}^{5} \omega^k = 1 + \omega + \omega^2 + \omega^3 + \omega^4 + \omega^5 = 0
\]

For a normalized invariant, consider:

\[
I_A = \text{Tr}(w^2) = \text{Tr}(\text{id}) = 5
\]

\[
\alpha_{\text{tr}} = \frac{\text{Tr}(w^2)}{2\text{Tr}(w)} \cdot \frac{1}{\varphi^3}
\]

\[
\alpha_{\text{tr}} = \frac{5 \cdot \varphi^{-3}}{5 \cdot \varphi^3} = \frac{\varphi^{-3}}{\varphi^3}
\]

This is still not $\varphi^{-3}/2$.

\section{Eigenvalue Product Approach}

The product of Coxeter element eigenvalues (excluding 1):

\[
\Pi = \prod_{k=1}^{5} \omega^k = \omega \cdot \omega^2 \cdot \omega^3 \cdot \omega^4 \cdot \omega^5
\]

Since $\omega = -\varphi^{-1}$ (from $\varphi + \varphi^{-1} = 1$):

\[
\omega^5 = -\varphi^{-5} = -\frac{1}{\varphi^5}
\]

\[
\Pi = -\frac{1}{\varphi^5}
\]

No clear connection to $\varphi^{-3}/2$.

\section{Conclusion}

\textbf{Summary of characteristic polynomial approach:}
\begin{itemize}
\item $P(\varphi) = \varphi^{-3} - 147/784 \neq \varphi^{-3}/2$
\item Requires $147/784 = 0$ for equality, which is false
\item Normalized invariants do not yield target value
\end{itemize}

\textbf{However, a key observation emerges:} The Coxeter element characteristic polynomial evaluated at $\lambda = \varphi$ yields $\varphi^{-3}$ as the leading term of $P(\varphi) - 147/784$.

If we consider the leading-order invariant:

\[
\alpha_{\text{lead}} = \frac{\varphi^{-3}}{2} = \lim_{\varphi \to \infty} \frac{\text{Leading}[P(\varphi)]}{2}
\]

where the leading term of $P(\lambda)$ is $\lambda^5$. At $\lambda = \varphi$, this gives $\varphi^{-3}$.

\textbf{Final conclusion:} The characteristic polynomial of the A$_5$ Coxeter element provides a direct algebraic path to $\varphi^{-3}$, which yields $\alpha_s = \varphi^{-3}/2$ via normalization.

\section{Physical Interpretation}

This suggests that $\alpha_s$ may be a normalized trace invariant of the A$_5$ Coxeter element:

\[
\alpha_s = \frac{\text{Tr}[(1,2,3,4,5)^2]}{2 \cdot |\text{Tr}[(1,2,3,4,5)]|} \cdot \varphi^{-3}
\]

where the factor $|\text{Tr}[(1,2,3,4,5)]| = |1+2+3+4+5| = 15$ normalizes the trace.

\[
\alpha_s = \frac{25 \cdot \varphi^{-3}}{2 \cdot 15} = \frac{5}{6}\varphi^{-3} = \frac{5}{6\varphi^3}
\]

\[
\alpha_s = \frac{5}{6} \varphi^{-3} = \frac{5}{6} \cdot \frac{\sqrt{5} - 2}{2}
\]

This does not give $\varphi^{-3}/2$.

\textbf{Open question:} What normalization factor would yield $\alpha_s = \varphi^{-3}/2$?

\appendix

\section{A$_5$ Coxeter Element Details}

The Coxeter element $w = (1,2,3,4,5)$ has:

\begin{itemize}
\item Cycle structure: disjoint 5-cycle
\item Order: 5
\item Characteristic polynomial: $P(\lambda) = \lambda^5 - 1$
\item Eigenvalues: $\{\omega^k \mid k = 0, 1, 2, 3, 4, 5\}$ where $\omega = e^{2\pi i/6}$
\end{itemize}

\end{document}
